% Chapter 2, Figure 5: angular acceleration of a flywheel (scan: figures/ch2/fig05.png).
% The orthographic view of Figure 3 (elevation 32 deg, azimuth 45 deg; TikZ x right and away, y left and
% away, z up). The wheel turns on the y axis, its front face (the plane y = 0) towards the viewer, a short
% hub standing out of the face; the wheel is opaque, so the axis is hidden behind it.
% The moment M and the angular velocity omega(t) (curved arrows on the face) turn the wheel clockwise on the
% page, a positive rotation about +y (x towards -z). The fixed reference line (dashed) is horizontal, along
% x, in the plane of the face; the line fixed on the wheel has turned from it through alpha(t) in the same
% sense (the 1973 arc has one arrowhead, at the wheel's line). alpha is drawn at 25 degrees (the scan: 24).
% The formulas are those of chapters/ch2-intro-sec1.tex:310-321, set beneath the wheel as printed.
\documentclass[11pt]{standalone}
\usepackage{tamrfig}
\begin{document}
\begin{tikzpicture}[x={(1.15cm,0.61cm)}, y={(-1.15cm,0.61cm)}, z={(0cm,1.38cm)}]
  \def\phis{48.5}                      % silhouette angle of a body along y in this view
  \def\R{1}\def\t{0.5}                 % wheel: radius, thickness (back face at y = t)
  \def\rh{0.07}\def\lh{0.5}            % hub: radius, length out of the face
  \def\al{25}                          % alpha(t) as drawn
  \def\Lr{2.15}                        % length of the two radial lines
  % axis: the part behind the wheel is hidden by its fill
  \draw[centerline] (0,2.3,0) -- (0,0,0);
  % the wheel: back rim (visible half), silhouettes, front face
  \draw[outline, fill=white] ({-\R*cos(\phis)},0,{-\R*sin(\phis)})
    -- plot[domain=0:180, samples=61] ({\R*cos(\phis+180-\x)}, \t, {\R*sin(\phis+180-\x)})
    -- ({\R*cos(\phis)},0,{\R*sin(\phis)}) -- cycle;
  \draw[outline, fill=white, canvas is xz plane at y=0] (0,0) circle[radius=\R];
  \begin{scope}[canvas is xz plane at y=0]
    % fixed horizontal reference line and the line on the wheel, turned through alpha(t)
    \draw[guide] (0:\rh) -- (0:\Lr);
    \draw[outline] (-\al:\rh) -- (-\al:\Lr);
    \draw[draw=ink, line width=0.5pt, -{Stealth[length=4pt,width=2.6pt]}]
      (0:1.85) arc[start angle=0, end angle=-\al, radius=1.85];
    % applied moment and angular velocity: clockwise on the page
    \draw[thin vec] (-58:0.5) arc[start angle=-58, end angle=-342, radius=0.5];
    \draw[thin vec] (-45:0.76) arc[start angle=-45, end angle=-343, radius=0.76];
  \end{scope}
  \node[right=1pt] at ({1.85*cos(\al/2)},0,{-1.85*sin(\al/2)}) {$\alpha(t)$};
  \node at ({0.3*cos(42)},0,{0.3*sin(42)}) {$M$};
  \node[anchor=south west, inner sep=1pt] (w) at ({1.12*cos(30)},0,{1.12*sin(30)}) {$\omega(t)$};
  \draw[leader, shorten >=1pt] (w.south west) -- ({0.76*cos(38)},0,{0.76*sin(38)});
  % hub: a short cylinder standing out of the face towards the viewer
  \draw[outline, fill=white] ({-\rh*cos(\phis)},0,{-\rh*sin(\phis)}) -- ({-\rh*cos(\phis)},-\lh,{-\rh*sin(\phis)})
    -- ({\rh*cos(\phis)},-\lh,{\rh*sin(\phis)}) -- ({\rh*cos(\phis)},0,{\rh*sin(\phis)});
  \draw[outline, canvas is xz plane at y=0] (\phis:\rh) arc[start angle=\phis, end angle=\phis+180, radius=\rh];
  \draw[outline, fill=white, canvas is xz plane at y=-\lh] (0,0) circle[radius=\rh];
  % axis through the hub and out in front
  \draw[centerline] (0,0,0) -- (0,-2.9,0);
  % the formulas, beneath the wheel
  \node[anchor=north west, inner sep=0pt, font=\small] at (-2.55cm,-1.6cm) {\setlength{\jot}{1pt}$\begin{aligned}
    \omega &= \gamma t\\
    &= \frac{M}{I}\,t\\
    \alpha &= \frac{\gamma}{2}\,t^{2}\\
    &= \frac{M}{2I}\,t^{2}
  \end{aligned}$};
\end{tikzpicture}
\end{document}
